\documentclass[10pt,a4paper]{amsart}
\usepackage[margin=19mm]{geometry}
\usepackage{amsmath,amssymb,mathtools}
\usepackage[scr=rsfs,cal=euler]{mathalfa}
\usepackage{xfrac}
\usepackage[unicode,hidelinks,bookmarks=false]{hyperref}
\newcommand{\ie}{i.e.\ }
\newcommand{\cf}{cf.\ }
\newcommand{\resp}{respectively\ }
\newcommand{\Subsection}{\subsection}
\providecommand{\nobreakdash}{\nobreak\hbox{-}}
\setlength{\emergencystretch}{2em}
\multlinegap0pt
\usepackage{mdframed}
\usepackage{mdframed}
\usepackage{mdframed}
\usepackage{mdframed}
\usepackage{mathtools}
\usepackage{amssymb}
\usepackage{bm}
\usepackage{tcolorbox}
\usepackage[shortlabels]{enumitem}
\usepackage{mdframed}
\newtheorem{theorem}{Theorem}
\DeclareMathOperator{\Gal}{Gal}
\newcommand{\cM}{\mathcal{M}}
\newcommand{\cQ}{\mathcal{Q}}
\title{Statistics of the Genus Number of $S_3 \times C_q$ and $D_4$-fields}
\author{ANUP B. DIXIT, Sunil Kumar Pasupulati}
\date{Source: arXiv:2605.04792v2 (CC BY 4.0)}
\begin{document}
\maketitle

\noindent\emph{Source and licence.} Statistics of the Genus Number of $S_3 \times C_q$ and $D_4$-fields. Version arXiv:2605.04792v2 (CC BY 4.0), distributed by arXiv under the Creative Commons Attribution 4.0 International licence, http://creativecommons.org/licenses/by/4.0/, as recorded in the arXiv licence metadata for this exact version and shown on the abstract page https://arxiv.org/abs/2605.04792v2. Full licence notice: \texttt{licenses/arxiv-2605.04792-CC-BY-4.0.txt}.\par\smallskip

\noindent\textit{Editorial scope.} Counted unit: the article's theorem S\_3-C\_q-theorem (source0.tex lines 278--283). The article's complete original proof of it is delivered with it (source0.tex lines 1014--1104, one \texttt{proof} environment of 91 lines, which the article places away from the statement). No mathematics is written, repaired or re-derived: the statement and the proof are the article's own lines, in the article's own notation, and no retained line is substituted or re-flowed.  Retained uncounted as the source-local context the proof needs: lines 1002--1006.  Nothing was omitted from any retained span, so the package carries no omission record.  No retained line contains a citation, so the wrapper carries no bibliography and no bibliography entry.  No retained line contains a figure environment, an image-inclusion command or drawing code, so no image is delivered and the package declares no asset dependency.  Editorial formatting only, in addition: an AMS \texttt{amsart} preamble that restates the article's own declarations and macros used by the retained lines, namely the environments \texttt{theorem} on separate counters, together with the standard packages those lines need.  The retained environments print in this document as Theorem 1 in order of appearance, whereas the article prints the same items with the numbering of its own full document.  The complete native source of the article as distributed in the arXiv e-print for this exact version is bundled unchanged as \texttt{sources/199896/source0.tex}.
\begin{equation}\label{lemma:count-genus-number-S_3}
    N_3^{\ell}(X)
    = \frac{\alpha_{\ell}}{\zeta(2)}\,X
      + O\!\left(X^{16/17+\varepsilon}\right),
\end{equation}

\begin{theorem}\label{S_3-C_q-theorem}
    For any prime $q\geq 5$ and $k,l\geq 0$, there exists a constant $C_q(k,l)>0$, explicitly described in \eqref{implied-constant-1}, such that
    \begin{equation*}
        |\cQ^q_{k,\ell}(X)| = C_q(k,\ell) \, X^{1/q} + O\left(X^{16/17q}\right).
    \end{equation*}
\end{theorem}

\begin{proof}[Proof of Theorem \ref{S_3-C_q-theorem}]
Since $[K:\mathbb{Q}]=3$ and $[F:\mathbb{Q}]=q$ are coprime, the fields $K$ and
$F$ are linearly disjoint and satisfy
\[
    \mathcal{D}_{KF}
    = \mathcal{D}_{K}^{\,q}\, \mathcal{D}_{F}^{\,3}.
\]
Thus the condition $|\mathcal{D}_{KF}|\le X$ is equivalent to
\[
    |\mathcal{D}_K|^{\,q}\,|\mathcal{D}_F|^{\,3} \le X.
\]
Therefore,
\begin{align*}
    |\cQ^q_{k,l}(X)|
    &= 
    \#\bigl\{\, L \in \cM(S_3\times C_q, 3q, X) : \ g_L = 3^{\ell}q^{k} \,\bigr\} \\
    &=
    \sum_{\substack{
        F:\, \Gal(F/\mathbb{Q})\simeq C_q\\
        g_F = q^{k},\ |\mathcal{D}_F|^{3}\le X}}
    \
    \sum_{\substack{
        K\ \text{cubic},\ \Gal(\widetilde{K}/\mathbb{Q})\simeq S_3\\
        g_K = 3^{\ell},\ |\mathcal{D}_K|^{q}\le X/|\mathcal{D}_F|^{3}}}
    1.
\end{align*}
Using \eqref{lemma:count-genus-number-S_3}, we obtain
\begin{align}\label{eqn-a-1}
    |\cQ^q_{k,l}(X)|
    &= 
    \sum_{\substack{
        F:\, \Gal(F/\mathbb{Q})\simeq C_q\\
        g_F = q^{k},\ |\mathcal{D}_F|\le X^{1/3}}}
    N_3^{\ell}\!\left( \left(\frac{X}{\mathcal{D}_F^3}\right)^{1/q} \right) \nonumber\\[6pt]
    &=
    \sum_{\substack{
        F:\, \Gal(F/\mathbb{Q})\simeq C_q\\
        \Psi_q(\mathcal{D}_F)=k+1,\ |\mathcal{D}_F|\le X^{1/3}}}
    \left[
        \frac{\alpha_\ell}{\zeta(2)}
        \left(\frac{X}{\mathcal{D}_F^3}\right)^{1/q}
        + O\!\left(
            \left(\frac{X}{\mathcal{D}_F^3}\right)^{16/(17q)+\varepsilon}
        \right)
    \right].
\end{align}
We show that the sum
\begin{equation}\label{eqn-sum}
    \sum_{\substack{
        F:\, \Gal(F/\mathbb{Q})\simeq C_q\\
        \Psi_q(\mathcal{D}_F)=k+1}}
    \frac{1}{\mathcal{D}_F^{3/q}}
\end{equation}
is convergent. For each cyclic extension $F/\mathbb{Q}$ of degree $q$, one has
$\mathcal{D}_F = f^{\,q-1}$, where $f$ is the conductor of $F$. Let $N_f(q)$ be the number of cyclic extensions of degree $q$ and conductor $f$.
As $N_f(q)$ is bounded by the number of order $q$ subgroups of
$(\mathbb{Z}/f\mathbb{Z})^\times$. If $q^{r}$ is the largest power of $q$ dividing $\varphi(f)$, then one has 
\[
    N_f(q)
    \le \frac{q^{r}-1}{q-1}
    \le \varphi(f)
    \le f.
\]
Thus
\begin{equation*}
\sum_{\substack{
        F:\, \Gal(F/\mathbb{Q})\simeq C_q \\ \Psi_q(\mathcal{D}_F)=k+1}
        }
\frac{1}{\mathcal{D}_F^{3/q}}
\ll
\sum_{\substack{\omega(f)=k+1}}
\frac{N_f(q)}{f^{3(q-1)/q}}
\ll
\sum_{\substack{\omega(f)=k+1}}
\frac{1}{f^{2-\frac{3}{q}}} \ll 1. 
\end{equation*}
Therefore, from Equation \eqref{eqn-a-1}, we conclude that
\begin{equation*}
    |\cQ^q_{k,l}(X)| = C_q(k,\ell) X^{1/q} + O(X^{16/17q}),
\end{equation*}
where
\begin{equation}\label{implied-constant-1}
    C_q(k,\ell)
    =
    \frac{\alpha_{\ell}}{\zeta(2)}
    \sum_{\substack{
        F:\, \Gal(F/\mathbb{Q})\simeq C_q\\
        \Psi_q(\mathcal{D}_F)=k+1}}
    \frac{1}{\mathcal{D}_F^{3/q}}.
\end{equation}
\end{proof}

\end{document}
